\documentclass[11pt,a4paper]{scrartcl}
\usepackage[provide=*,lithuanian]{babel}
\usepackage{../1.\space Vienanariai\space ir\space daugianariai/manostilius_lt}
\usepackage[margin=20mm,headheight=15pt,headsep=7mm,footskip=10mm]{geometry}
\usepackage{multicol,booktabs,needspace}
\usetikzlibrary{arrows.meta,decorations.pathreplacing}
\hypersetup{hidelinks,pdftitle={Kvadratinio trinario reikšmių simetrija},pdfauthor={Andrius Berniukevičius}}
\title{Trinario reikšmių simetrija}
\titlehead{}
\setlength{\parskip}{4pt}
\setlist[enumerate]{leftmargin=*,itemsep=7pt,topsep=4pt}
\setlength{\columnsep}{9mm}
\newcommand{\pagetitle}[2]{\begin{center}{\Large\bfseries #1}\\[3pt]{\small #2}\end{center}\vspace{3mm}}
\newcommand{\pv}[2]{\begin{example}\label{#1}#2\end{example}}
\newcommand{\hint}[1]{{\small\itshape Jei kyla sunkumų, žr. #1 pavyzdžius.}\par}
\newcommand{\numberlinefig}[6]{%
\begin{center}\begin{tikzpicture}[x=1cm,y=1cm,>=Latex]
\draw[->,gray] (-.6,0)--(8.6,0);
\foreach \x/\t in {0/{#1},4/{#2},8/{#3}}{\draw (\x,-.09)--(\x,.09);\node[below=4pt] at (\x,0){$\t$};}
\fill[MidnightBlue] (0,0) circle(2pt) (8,0) circle(2pt);
\fill (4,0) circle(2pt);
\draw[<->,MidnightBlue] (0,.55)--node[above]{#4}(4,.55);
\draw[<->,MidnightBlue] (4,.55)--node[above]{#5}(8,.55);
\node[below=20pt,font=\small] at (4,0){#6};
\end{tikzpicture}\end{center}}
\begin{document}

% TEORIJA
\pagetitle{Kvadratinio trinario reikšmių simetrija}{Atstumai nuo ekstremumo taško · teorija ir pavyzdžiai}
\textbf{Pradinės žinios.} Pilnojo kvadrato išskyrimas; kvadratinio trinario ekstremumo taškas, minimumas ir maksimumas; skaičių tiesė.

\section{Atstumas nuo ekstremumo taško}
Imkime kvadratinį trinarį
\[P(x)=2(x-5)^2+3.\]
Jo minimumo taškas yra $x_{\min}=5$, o minimumas $P_{\min}=3$.
Skaičiai $3$ ir $7$ nuo $5$ nutolę po du vienetus, tik į skirtingas puses. Todėl skirtumai $3-5=-2$ ir $7-5=2$ yra priešingi, o jų kvadratai lygūs:
\[P(3)=2\cdot(-2)^2+3=11,\qquad P(7)=2\cdot2^2+3=11.\]
\numberlinefig{3}{5}{7}{2 vienetai}{2 vienetai}{$P(3)=P(7)$}
\begin{center}\begin{tabular}{ccc}\toprule
Kintamojo reikšmės & Atstumas nuo $5$ & Trinario reikšmė\\\midrule
$4$ ir $6$ & $1$ & $5$\\
$3$ ir $7$ & $2$ & $11$\\
$2$ ir $8$ & $3$ & $21$\\\bottomrule
\end{tabular}\end{center}
Trinario reikšmei svarbus atstumas nuo minimumo taško. Kurioje jo pusėje pasirinktas skaičius, įtakos neturi. Pagrįskime tai bendruoju atveju.

\begin{property}\label{prop:sym}
Jei dvi kintamojo reikšmės vienodu atstumu nutolusios nuo kvadratinio trinario ekstremumo taško, tai atitinkamos trinario reikšmės yra lygios.
\end{property}
\begin{proof}
Išskyrę pilnąjį kvadratą, užrašome $P(x)=a(x-h)^2+k$, kur $a\ne0$, o $h$ yra ekstremumo taškas. Jei atstumas nuo pasirinkto skaičiaus iki $h$ yra $d\ge0$, tai $x-h=d$ arba $x-h=-d$. Abiem atvejais $(x-h)^2=d^2$, todėl $P(x)=ad^2+k$. Taigi vienodiems atstumams atitinka vienodos trinario reikšmės.
\end{proof}
Trumpai ši savybė užrašoma taip:
\[\boxed{P(h-d)=P(h+d).}\]
Ji galioja tiek minimumo, tiek maksimumo atveju.

\newpage
\section{Kaip rasti tokią pačią trinario reikšmę?}
\pv{pv:positive}{Trinario ekstremumo taškas yra $5$, o $P(2)=9$. Raskite $P(8)$.}
\begin{solution}
Nuo $2$ iki $5$ yra trys vienetai. Nuo $5$ iki $8$ taip pat trys vienetai. Atstumai nuo ekstremumo taško vienodi, todėl $P(8)=P(2)=9$.
\end{solution}
\pv{pv:negative}{Trinario ekstremumo taškas yra $-2$. Raskite kitą kintamojo reikšmę, su kuria trinario reikšmė tokia pati kaip su $x=3$.}
\begin{solution}
Skaičius $3$ yra penkiais vienetais dešiniau už $-2$. Penkiais vienetais kairiau už $-2$ yra skaičius $-7$. Todėl $P(-7)=P(3)$.
\end{solution}
\numberlinefig{-7}{-2}{3}{5 vienetai}{5 vienetai}{$P(-7)=P(3)$}

\begin{property}\label{prop:reverse}
Jei su dviem skirtingomis kintamojo reikšmėmis gaunamos vienodos kvadratinio trinario reikšmės, tai ekstremumo taškas yra per vidurį tarp tų skaičių.
\end{property}
\begin{proof}
Tegu $P(x)=a(x-h)^2+k$, o $P(u)=P(v)$ ir $u\ne v$. Atėmę $k$ ir padaliję iš $a\ne0$, gauname $(u-h)^2=(v-h)^2$. Vadinasi, atstumų nuo $u$ ir $v$ iki $h$ kvadratai lygūs. Patys atstumai taip pat lygūs, nes jie neneigiami. Skirtingi skaičiai toje pačioje $h$ pusėje negali būti vienodai nuo jo nutolę. Todėl $u$ ir $v$ yra skirtingose $h$ pusėse, o $h$ yra jų vidurys.
\end{proof}
\pv{pv:midpoint}{Žinoma, kad $P(8)=P(14)$. Raskite ekstremumo tašką ir kitą skaičių, su kuriuo trinario reikšmė lygi $P(9)$.}
\begin{solution}
Tarp $8$ ir $14$ yra šeši vienetai; pusė šio atstumo yra trys vienetai. Ekstremumo taškas yra $8+3=11$. Skaičius $9$ yra dviem vienetais kairiau už $11$, o $13$ -- dviem vienetais dešiniau. Todėl $P(9)=P(13)$.
\end{solution}
\textbf{Pastaba.} Sąlyga, kad skaičiai skirtingi, būtina. Lygybė $P(8)=P(8)$ ekstremumo taško nenusako.

\newpage
\section{Vienodų reikšmių paieška ir reikšmių palyginimas}
\pv{pv:equation}{Trinario ekstremumo taškas yra $4$. Išspręskite lygtį $P(x)=P(1)$.}
\begin{solution}
Skaičius $1$ nuo $4$ nutolęs trimis vienetais į kairę. Tuo pačiu atstumu į dešinę yra $7$. Todėl tinka $x=1$ ir $x=7$. Kitų sprendinių nėra: pagal \ref{prop:reverse} savybę kiekvienas nuo $1$ skirtingas sprendinys turi būti simetriškas jam ekstremumo taško atžvilgiu. Atsakymas: $x=1$ arba $x=7$.
\end{solution}
Jei lygties dešinėje būtų $P(4)$, vienintelis sprendinys būtų $x=4$. Nuliniu atstumu nuo ekstremumo taško nutolęs tik pats ekstremumo taškas.

\begin{property}\label{prop:compare}
Kai trinaris turi minimumą, didesniam atstumui nuo minimumo taško atitinka didesnė trinario reikšmė. Kai trinaris turi maksimumą, didesniam atstumui nuo maksimumo taško atitinka mažesnė trinario reikšmė.
\end{property}
\begin{proof}
Su atstumu $d$ trinario reikšmė lygi $ad^2+k$. Didėjant neneigiamam atstumui, jo kvadratas didėja. Kai $a>0$, dauginant iš $a$ tvarka išlieka, o kai $a<0$ -- apsiverčia. Pridėjus $k$, tvarka nesikeičia.
\end{proof}
\pv{pv:min}{Trinario minimumo taškas yra $4$. Palyginkite $P(1)$ ir $P(6)$.}
\begin{solution}
Skaičius $1$ nuo $4$ nutolęs trimis vienetais, o $6$ -- dviem. Toliau nuo minimumo taško gaunama didesnė trinario reikšmė, todėl $P(1)>P(6)$.
\end{solution}
\pv{pv:max}{Žinoma, kad $P(2)=P(10)$, o trinaris turi maksimumą. Palyginkite $P(3)$ ir $P(8)$.}
\begin{solution}
Ekstremumo taškas yra $6$. Skaičius $3$ nuo jo nutolęs trimis vienetais, o $8$ -- dviem. Toliau nuo maksimumo taško gaunama mažesnė trinario reikšmė, todėl $P(3)<P(8)$.
\end{solution}
\textbf{Pastaba.} Vien lygybės $P(2)=P(10)$ nepakanka skirtingiems atstumams atitinkančioms reikšmėms palyginti. Reikia žinoti, ar trinaris turi minimumą, ar maksimumą.

\newpage
\section{Kaip sujungti kelias savybes?}
\pv{pv:restore}{Raskite kvadratinį trinarį, jei $P(1)=P(7)=5$, o $P(4)=-4$.}
\begin{solution}
Ekstremumo taškas yra per vidurį tarp $1$ ir $7$, taigi jis lygus $4$. Kadangi $P(4)=-4$, tai
\[P(x)=a(x-4)^2-4.\]
Skaičius $1$ nuo ekstremumo taško nutolęs trimis vienetais. Jo atstumo kvadratas lygus $9$, todėl $5=9a-4$. Gauname $a=1$. Vadinasi,
\[P(x)=(x-4)^2-4=x^2-8x+12.\]
\end{solution}
\pv{pv:distance-square}{Trinario minimumo taškas yra $5$, $P(5)=2$, o $P(6)=5$. Raskite $P(3)$.}
\begin{solution}
Vieno vieneto atstumu nuo minimumo taško trinario reikšmė už minimumą didesnė $5-2=3$. Dviejų vienetų atstumo kvadratas yra keturis kartus didesnis už vieno vieneto atstumo kvadratą. Iš pavidalo $ad^2+k$ matyti, kad skirtumas nuo minimumo taip pat keturis kartus didesnis. Todėl
\[P(3)=2+4\cdot3=14.\]
Dauginame skirtumą nuo minimumo, o ne pačią trinario reikšmę.
\end{solution}

\subsection*{Ką verta įsidėmėti?}
\begin{enumerate}
\item Vienodu atstumu nuo ekstremumo taško nutolusiems skaičiams atitinka vienodos trinario reikšmės.
\item Dvi skirtingos kintamojo reikšmės, su kuriomis trinaris įgyja tą pačią reikšmę, yra skirtingose ekstremumo taško pusėse. Jų vidurys yra ekstremumo taškas.
\item Lygindami reikšmes, pirmiausia raskite atstumus nuo ekstremumo taško, tada atsižvelkite į tai, ar trinaris turi minimumą, ar maksimumą.
\end{enumerate}
\textbf{Formulių apibendrinimas.} Jei ekstremumo taškas yra $h$, tai skaičiui $u$ simetriškas skaičius yra $2h-u$. Jei $P(u)=P(v)$ ir $u\ne v$, tai $h=\frac{u+v}{2}$. Šios formulės užrašo jau aptartus samprotavimus apie atstumus ir vidurį.


% UZDAVINIAI
\newpage
\pagetitle{Uždaviniai}{Visur $P(x)$ yra kvadratinis trinaris; $a$ -- jo pirmasis koeficientas.}
\textit{Samprotaukite per atstumus. Kai prašoma palyginti ar pagrįsti, nurodykite ekstremumo tašką ir reikalingus atstumus.}
\subsection*{A. Vienodi atstumai}
\hint{\ref{pv:positive} ir \ref{pv:negative}}
\begin{multicols}{2}\begin{enumerate}
\item Ekstremumo taškas yra $5$. Užbaikite: a) $P(3)=P(\ldots)$; b) $P(9)=P(\ldots)$; c) $P(-2)=P(\ldots)$.
\item Ekstremumo taškas yra $-3$. Užbaikite: a) $P(-5)=P(\ldots)$; b) $P(2)=P(\ldots)$; c) $P(0)=P(\ldots)$.
\item Ekstremumo taškas yra $4$, $P(-1)=12$. Raskite $P(9)$.
\item Ekstremumo taškas yra $-2$, $P(-6)=-7$. Raskite $P(2)$.
\item $P(x)=-3(x+2)^2+8$. Prie kiekvieno skaičiaus $0$, $3$, $-7$ parašykite kitą skaičių, su kuriuo trinario reikšmė tokia pati.
\item Ekstremumo taškas yra $\frac12$. Raskite skaičiui $-\frac32$ simetrišką skaičių.
\end{enumerate}\end{multicols}\subsection*{B. Vidurio radimas ir vienodų reikšmių paieška}\hint{\ref{pv:midpoint} ir \ref{pv:equation}}\begin{multicols}{2}\begin{enumerate}[start=7]
\item Raskite ekstremumo tašką: a) $P(1)=P(9)$; b) $P(-7)=P(3)$; c) $P(4)=P(5)$.
\item $P(-2)=P(8)$. Raskite $t\ne0$, su kuriuo $P(t)=P(0)$.
\item $P(3)=P(11)$. Kurios lygybės teisingos: a) $P(6)=P(8)$; b) $P(0)=P(14)$; c) $P(5)=P(10)$?
\item Ekstremumo taškas yra $3$. Išspręskite: a) $P(x)=P(1)$; b) $P(x)=P(8)$; c) $P(x)=P(3)$.
\item $P(-4)=P(6)$. Išspręskite $P(x)=P(3)$.
\item Ekstremumo taškas yra $-1$. Išspręskite $P(x)=P(2)$.
\end{enumerate}\end{multicols}\newpage\subsection*{C. Reikšmių palyginimas}\hint{\ref{pv:min} ir \ref{pv:max}}\begin{multicols}{2}\begin{enumerate}[start=13]
\item Minimumo taškas yra $5$. Palyginkite: a) $P(4)$ ir $P(8)$; b) $P(2)$ ir $P(8)$; c) $P(5)$ ir $P(6)$.
\item Maksimumo taškas yra $-2$. Palyginkite: a) $P(-5)$ ir $P(0)$; b) $P(-4)$ ir $P(0)$; c) $P(-2)$ ir $P(1)$.
\item $P(2)=P(10)$, $a<0$. Išdėstykite didėjimo tvarka $P(6)$, $P(5)$, $P(2)$.
\item $P(-3)=P(5)$, $a>0$. Palyginkite $P(-1)$ ir $P(4)$.
\end{enumerate}\end{multicols}\subsection*{D. Ar pakanka duomenų? Ar tai įmanoma?}\begin{multicols}{2}\begin{enumerate}[start=17]
\item Ar iš $P(8)=P(12)$ galima nustatyti, kuri reikšmė didesnė: $P(9)$ ar $P(13)$? Kokios informacijos pakaktų?
\item Ar iš $P(4)=P(4)$ galima spręsti, kad ekstremumo taškas yra $4$? Pateikite kontrpavyzdį.
\item Ar gali kvadratinis trinaris tenkinti abi sąlygas: $P(1)=P(7)$ ir $P(2)=P(8)$?
\item Ar kvadratinis trinaris gali įgyti tą pačią reikšmę su trimis skirtingais skaičiais? Pagrįskite.
\end{enumerate}\end{multicols}\newpage\subsection*{E. Kelių savybių taikymas}\hint{\ref{pv:restore} ir \ref{pv:distance-square}}\begin{enumerate}[start=21]
\item $P(0)=P(6)=10$, $P(3)=1$. Raskite $P(x)$.
\item $P(-1)=P(5)=4$, $P(2)=13$. Raskite $P(x)$.
\item $P(2)=P(8)$, $P(4)=7$, $P(5)=5$. Raskite $P(6)$ ir $P(3)$.
\item Minimumo taškas $-1$, $P(-1)=-3$, $P(0)=2$. Raskite $P(-3)$.
\end{enumerate}\subsection*{F. Savarankiškas samprotavimas}\begin{enumerate}[start=25]
\item $P(1)=P(9)$. Įrodykite, kad $P(2)+P(8)=2P(2)$.
\item $P(2)=P(10)$, $P(3)+P(9)=18$. Raskite $P(3)$.
\item $P(1)=P(7)$ ir $P(2)>P(3)$. Ar trinaris turi minimumą, ar maksimumą?
\item $P(u)=P(v)$, $u<v$. Įrodykite, kad kiekvienam skaičiui $t$ galioja $P(u+t)=P(v-t)$.
\end{enumerate}
\subsection*{Trumpas pasitikrinimas vėliau}
Uždenkite atsakymus. Dar kartą išspręskite 2, 8, 14 ir 27 uždavinius. Kiekviename sprendime aiškiai pasakykite, kokį vaidmenį atlieka atstumas nuo ekstremumo taško.


% ATSAKYMAI
\newpage\pagetitle{Atsakymai}{Trumpi atsakymai į 1--28 uždavinius}\begin{multicols}{2}\begin{enumerate}
\item a) $7$; b) $1$; c) $12$.
\item a) $-1$; b) $-8$; c) $-6$.
\item $12$.
\item $-7$.
\item Atitinkamai $-4$, $-7$, $3$.
\item $\frac52$.
\item a) $5$; b) $-2$; c) $\frac92$.
\item $t=6$.
\item a) ir b); c) neteisinga.
\item a) $1$, $5$; b) $-2$, $8$; c) $3$.
\item $x=-1$ arba $x=3$.
\item $x=-4$ arba $x=2$.
\item a) $<$; b) $=$; c) $<$.
\item a) $<$; b) $=$; c) $>$.
\item $P(2)<P(5)<P(6)$.
\item $P(-1)<P(4)$.
\item Nepakanka; pakaktų $a$ ženklo.
\item Ne; pvz., $P(x)=x^2+1$.
\item Negali.
\item Negali.
\item $P(x)=(x-3)^2+1$.
\item $P(x)=-(x-2)^2+13$.
\item $P(6)=7$, $P(3)=13$.
\item $17$.
\item Lygybė teisinga, nes $P(2)=P(8)$.
\item $9$.
\item Minimumą.
\item Teiginys teisingas.
\end{enumerate}\end{multicols}


% SPRENDIMAI
\newpage\pagetitle{Mokytojo sprendimai}{Atstumais pagrįsti visų uždavinių sprendimai}\begin{enumerate}
\item Atstumai nuo $5$ yra atitinkamai $2$, $4$, $7$. Į kitą pusę tiek pat nutolę skaičiai yra $7$, $1$, $12$.
\item Atstumai nuo $-3$ yra $2$, $5$, $3$. Kitoje pusėje gauname $-1$, $-8$, $-6$.
\item Abu skaičiai nuo $4$ nutolę po $5$ vienetus, todėl $P(9)=P(-1)=12$.
\item Atstumai nuo $-2$ lygūs $4$, todėl $P(2)=P(-6)=-7$.
\item Ekstremumo taškas $-2$. Atstumai yra $2$, $5$, $5$; kitoje pusėje gauname $-4$, $-7$, $3$.
\item Nuo $-\frac32$ iki $\frac12$ yra $2$ vienetai. Pridėję $2$ prie $\frac12$, gauname $\frac52$.
\item Randame atstumų puses: a) $1+8:2=5$; b) $-7+10:2=-2$; c) $4+1:2=\frac92$.
\item Ekstremumo taškas $3$. Nulis yra $3$ vienetais kairiau, o $6$ -- tiek pat dešiniau.
\item Ekstremumo taškas $7$. a) atstumai $1$ ir $1$; b) $7$ ir $7$; c) $2$ ir $3$. Skirtingi atstumai duoda skirtingas reikšmes pagal \ref{prop:compare} savybę.
\item a) Atstumas $2$, todėl tinka $1$ ir $5$. b) Atstumas $5$, todėl tinka $-2$ ir $8$. c) Atstumas $0$, todėl tinka tik $3$.
\item Ekstremumo taškas $1$. Nuo $3$ iki jo yra $2$ vienetai. Kitoje pusėje yra $-1$.
\item Skaičius $2$ nuo $-1$ nutolęs $3$ vienetais. Tiek pat į kairę nutolęs $-4$.
\item Atstumai nuo minimumo taško: a) $1<3$; b) $3=3$; c) $0<1$. Didesnis atstumas duoda didesnę reikšmę.
\item Atstumai: a) $3>2$; b) $2=2$; c) $0<3$. Toliau nuo maksimumo taško reikšmė mažesnė.
\item Maksimumo taškas $6$. Atstumai $4$, $1$, $0$. Kadangi $a<0$, didžiausiam atstumui atitinka mažiausia reikšmė.
\item Minimumo taškas $1$. Atstumai iki $-1$ ir $4$ yra $2$ ir $3$. Pirmasis mažesnis.
\item Ekstremumo taškas $10$, atstumai $1$ ir $3$. Kai $a>0$, $P(9)<P(13)$; kai $a<0$, $P(9)>P(13)$. Abu atvejai suderinami su duota lygybe.
\item Lygybė su tuo pačiu skaičiumi teisinga visada. Trinario $x^2+1$ ekstremumo taškas $0$, nors $P(4)=P(4)$.
\item Pirmoji sąlyga duoda ekstremumo tašką $4$, antroji -- $5$. Tačiau pavidale $a(x-h)^2+k$ ekstremumo taškas yra vienintelis: tik kai $x=h$ atstumo kvadratas lygus nuliui.
\item Tai pačiai trinario reikšmei turi atitikti tas pats atstumas nuo ekstremumo taško. Duotu teigiamu atstumu skaičių tiesėje yra tik du skaičiai, o nuliniu -- vienas.
\item Ekstremumo taškas $3$, ekstremumas $1$. Todėl $P(x)=a(x-3)^2+1$. Atstumas nuo $0$ iki $3$ yra $3$, tad $10=9a+1$ ir $a=1$.
\item Ekstremumo taškas $2$, ekstremumas $13$. Atstumas iki $-1$ yra $3$, todėl $4=9a+13$ ir $a=-1$.
\newpage
\item Ekstremumo taškas $5$. Skaičiai $4$ ir $6$ nuo jo nutolę po $1$, todėl $P(6)=7$. Vieno vieneto atstumu skirtumas nuo ekstremumo lygus $7-5=2$. Dviejų vienetų atstumu skirtumas keturis kartus didesnis, todėl $P(3)=5+4\cdot2=13$.
\item Vieno vieneto atstumu skirtumas nuo minimumo yra $2-(-3)=5$. Skaičius $-3$ nutolęs $2$ vienetais, todėl $P(-3)=-3+4\cdot5=17$.
\item Ekstremumo taškas $5$. Skaičiai $2$ ir $8$ nuo jo nutolę po $3$, todėl jų reikšmės lygios. Jų suma lygi dvigubai vienai iš jų.
\item Ekstremumo taškas $6$. Skaičiai $3$ ir $9$ nuo jo nutolę po $3$, todėl jų reikšmės lygios. Kiekviena lygi pusei sumos: $18:2=9$.
\item Ekstremumo taškas $4$. Skaičius $2$ nutolęs $2$ vienetais, o $3$ -- vienu. Tolimesniam skaičiui atitinka didesnė reikšmė. Maksimumo atveju būtų priešingai, todėl trinaris turi minimumą.
\item Ekstremumo taškas yra $h=(u+v)/2$. Skirtumai nuo jo yra priešingi: $(u+t)-h=-(v-t-h)$. Todėl atstumai nuo $h$ vienodi, net jei pasirinkti skaičiai peržengia ekstremumo tašką. Pagal simetrijos savybę jų trinario reikšmės lygios.
\end{enumerate}

\end{document}